\section{Contexto}

\paragraph{}
O Laboratório de Segurança disponibiliza um ambiente controlado destinado à
realização de atividades práticas de segurança ofensiva e testes de penetração.
O ambiente permite a execução de procedimentos de reconhecimento, análise de
vulnerabilidades e validação de falhas de segurança em sistemas propositalmente
disponibilizados para fins educacionais.

\paragraph{}
Durante a atividade de reconhecimento foi analisada a aplicação web
\texttt{http://testasp.vulnweb.com}, utilizada como alvo para a realização da
varredura automatizada de vulnerabilidades com a ferramenta OWASP ZAP. A
aplicação apresentou diversas páginas, rotas e parâmetros acessíveis por meio
de requisições HTTP, permitindo o mapeamento de funcionalidades e superfícies
de ataque disponíveis para análise.

\paragraph{}
A avaliação identificou componentes típicos de aplicações web dinâmicas,
incluindo mecanismos de autenticação, formulários de entrada de dados,
funcionalidades de pesquisa e páginas com parâmetros manipuláveis pelo usuário.
Também foram observados cabeçalhos HTTP, mensagens de erro e informações do
servidor que auxiliaram na caracterização do ambiente e de seus controles de
segurança.

\paragraph{}
Os resultados da varredura apontaram vulnerabilidades relevantes, entre elas
\textit{SQL Injection}, \textit{Cross Site Scripting (Reflected)}, \textit{Path
Traversal}, \textit{External Redirect}, ausência de proteção contra
\textit{Cross-Site Request Forgery (CSRF)} e falta de cabeçalhos de segurança
como \textit{Content Security Policy (CSP)} e mecanismos de proteção contra
\textit{Clickjacking}. Essas evidências indicam que a principal área de risco
do ambiente encontra-se na camada de aplicação web.

\paragraph{}
Por se tratar de um laboratório destinado ao treinamento e capacitação em
segurança da informação, todas as atividades são executadas em ambiente
controlado e isolado de sistemas de produção, permitindo a realização de testes
de forma segura e sem impactos a infraestruturas reais.
